% SOURCE_AUTHORITY: sga5_fr_workpass.tex, lines 7088--7130 (LF-normalized unit).
% SOURCE_SHA256: 791F4EFFC5E02832D5D77ED03518C8156D6F07E4C8238B03545DB93D883FBB28
\subsection*{3. Sistemas projetivos $J$-ádicos e AR-$J$-ádicos}
Sejam $A$ um anel comutativo com unidade, $J$ um ideal de $A$ e $(\cC,\rho)$ uma $A$-categoria abeliana. Supõe-se ou bem que $J$ é de tipo finito, ou bem que $\cC$ possui somas diretas infinitas. As notações são as mesmas que no número anterior. Em particular, os símbolos $P$ e $P_{AR}$ designam respectivamente as categorias $\Hom(\NN^\circ,\cC)$ e $\Hom_{AR}(\NN^\circ,\cC)$.

\subsubsection*{3.1. Sistemas projetivos $J$-ádicos}
\begin{statement}{Definição 3.1.1}
Seja $X$ um objeto de $P$. Diz-se que $X$ é $J$-ádico se verifica as seguintes condições:
\begin{enumerate}[label=(\roman*)]
\item para todo inteiro $n\ge0$, $J^{n+1}X_n=0$;
\item Para cada par $(m,n)$ de inteiros com $m\ge n\ge0$, o morfismo
\[
A/J^{n+1}\otimes_A X_m\longrightarrow X_n
\]
deduzido do morfismo de transição $X_m\to X_n$ é um isomorfismo.
\end{enumerate}
\end{statement}
Denotar-se-á por $J$-ad$(\cC)$ a subcategoria plena de $P$ engendrada pelos sistemas projetivos $J$-ádicos.

\subsubsection*{Observação 3.1.2}
Sejam $I$ e $J$ dois ideais de $A$ e suponhamos que $I$ e $J$ sejam de tipo finito, ou que $\cC$ admite somas diretas infinitas. Então pode ser demonstrado que, se $I$ e $J$ definem a mesma topologia sobre $A$, as categorias $I$-ad$(\cC)$ e $J$-ad$(\cC)$ são equivalentes.

\begin{statement}{Proposição 3.1.3 (Estabilidade por sequências exatas)}
Seja
\[
0\longrightarrow X\stackrel{u}{\longrightarrow}Y\stackrel{v}{\longrightarrow}Z\longrightarrow0
\]
uma sequência exata de $P$.
\begin{enumerate}[label=(\roman*)]
\item Se $X$ é estrito e $Y$ é $J$-ádico, então $Z$ é $J$-ádico.
\item Se $Y$ é estrito e $Z$ é $J$-ádico, então $X$ é estrito.
\item Se $X$ e $Z$ são $J$-ádicos, $Y$ é $J$-ádico.
\end{enumerate}
\end{statement}

\begin{prooftext}[Demonstração]
Se $m\ge n$, é claro que o diagrama $(T)$ que aparece a seguir, no qual as flechas verticais designam os morfismos de transição, é comutativo e é exato:
\[
\begin{tikzcd}[column sep=large,row sep=large]
{} & A/J^{n+1}\otimes_A X_m \arrow[r,"\id\otimes u_m"] \arrow[d,"f_{nm}"'] & A/J^{n+1}\otimes_A Y_m \arrow[r,"\id\otimes v_m"] \arrow[d,"g_{nm}"] & A/J^{n+1}\otimes_A Z_m \arrow[r] \arrow[d,"h_{nm}"] & 0\\
0 \arrow[r] & X_n \arrow[r] & Y_n \arrow[r] & Z_n \arrow[r] & 0 .
\end{tikzcd}
\]
No caso (i), $f_{nm}$ é um epimorfismo e $g_{nm}$ é um isomorfismo. No caso (ii), $g_{nm}$ é um epimorfismo e $h_{nm}$ um isomorfismo. No caso (iii), $f_{nm}$ e $h_{nm}$ são isomorfismos. Em todos os casos, o lema da serpente aplicado ao diagrama $(T)$ permite concluir.
\end{prooftext}
